% sec2_5.tex — Section 2.5: Estimating E(epsilon^2 x_i x_i') Consistently

The theory developed so far presumes that there is available a consistent estimator,
$\widehat{\mathbf{S}}$, of $\mathbf{S} = \E(\mathbf{g}_i \mathbf{g}_i') \equiv \E(\varepsilon_i^2 \mathbf{x}_i \mathbf{x}_i')$
to be used to calculate the estimated asymptotic
variance, $\widehat{\avar}(\mathbf{b})$.  This section explains how to obtain $\widehat{\mathbf{S}}$
from the sample $(\mathbf{y}, \mathbf{X})$.

%% ─── Using Residuals for the Errors ───────────────────────────────────────────
\subsection*{Using Residuals for the Errors}

If the error were observable, then the sample mean of $\varepsilon_i^2 \mathbf{x}_i \mathbf{x}_i'$
is obviously consistent by ergodic stationarity.  But we do not observe the error term,
and the substitution of some consistent estimate of it results in
\begin{equation}
  \widehat{\mathbf{S}} \equiv \frac{1}{n} \sum_{i=1}^{n} \hat{\varepsilon}_i^2\,
    \mathbf{x}_i \mathbf{x}_i',
  \label{eq:2.5.1}
\end{equation}
where $\hat{\varepsilon}_i \equiv y_i - \mathbf{x}_i' \widetilde{\bm{\beta}}$, and
$\widetilde{\bm{\beta}}$ is some consistent estimator of $\bm{\beta}$.
(Although the obvious candidate for the consistent estimator
$\widetilde{\bm{\beta}}$ is the OLS estimator $\mathbf{b}$, we use
$\widetilde{\bm{\beta}}$ rather than $\mathbf{b}$ here, in order to make the point
that the results of this section hold for any consistent estimator.)
For this estimator to be consistent for $\mathbf{S}$, we need to make a fourth-moment
assumption about the regressors.

\begin{assumption}[finite fourth moments for regressors]
\label{ass:2.6}
$\E[(x_{ik} x_{ij})^2]$ \emph{exists and is finite for all} $k$, $j$ $(= 1, 2, \ldots, K)$.
\end{assumption}

\begin{proposition}[consistent estimation of $\mathbf{S}$]
\label{prop:2.4}
\emph{Suppose the coefficient estimate $\widetilde{\bm{\beta}}$ used for calculating the
residual $\hat{\varepsilon}_i$ for $\widehat{\mathbf{S}}$ in \eqref{eq:2.5.1} is consistent,
and suppose $\mathbf{S} = \E(\mathbf{g}_i \mathbf{g}_i')$ exists and is finite.  Then,
under Assumptions~2.1, 2.2, and 2.6, $\widehat{\mathbf{S}}$ given in \eqref{eq:2.5.1} is
consistent for $\mathbf{S}$.}
\end{proposition}

To indicate why the fourth-moment assumption is needed for the regressors, we provide
a sketch of the proof for the special case of $K = 1$ (only one regressor).  So
$\mathbf{x}_i$ is now a scalar $x_i$, $\mathbf{g}_i$ is a scalar $g_i = x_i \varepsilon_i$,
and (2.3.10) (with $\mathbf{b} = \widetilde{\bm{\beta}}$ and $e_i = \hat{\varepsilon}_i$)
simplifies to
\begin{equation}
  \hat{\varepsilon}_i^2 = \varepsilon_i^2 - 2(\widetilde{\beta} - \beta) x_i \varepsilon_i
    + (\widetilde{\beta} - \beta)^2 x_i^2.
  \label{eq:2.5.2}
\end{equation}
By multiplying both sides by $x_i^2$ and summing over $i$,
\begin{equation}
  \frac{1}{n} \sum_{i=1}^{n} \hat{\varepsilon}_i^2 x_i^2
    - \frac{1}{n} \sum_{i=1}^{n} \varepsilon_i^2 x_i^2
    = -2(\widetilde{\beta} - \beta) \frac{1}{n} \sum_{i=1}^{n} \varepsilon_i x_i^3
    + (\widetilde{\beta} - \beta)^2 \frac{1}{n} \sum_{i=1}^{n} x_i^4.
  \label{eq:2.5.3}
\end{equation}
Now we can see why the finite fourth-moment assumption on $x_i$ is required: if the
fourth moment $\E(x_i^4)$ is finite, then by ergodic stationarity the sample average of
$x_i^4$ converges in probability to some finite number, so that the last term in
\eqref{eq:2.5.3} vanishes (converges to 0 in probability) if $\widetilde{\beta}$ is
consistent for $\beta$.  It can also be shown (see Analytical Exercise~4 for proof) that,
by combining the same fourth-moment assumption and the fourth-moment assumption that
$\E(\mathbf{g}_i \mathbf{g}_i') \;(= \E(\varepsilon_i^2 \mathbf{x}_i \mathbf{x}_i'))$ is
finite, the sample average of $x_i^3 \varepsilon_i$ converges in probability to some
finite number, so that the other term on the RHS of \eqref{eq:2.5.3}, too, vanishes.

According to Proposition~2.1(a), the assumptions made in Proposition~2.3 are sufficient
to guarantee that $\mathbf{b}$ is consistent, so we can set
$\mathbf{b} = \widetilde{\bm{\beta}}$ in \eqref{eq:2.5.1} and use the OLS residual to
calculate $\widehat{\mathbf{S}}$.  Also, the assumption made in Proposition~2.4 that
$\E(\mathbf{g}_i \mathbf{g}_i')$ is finite is part of Assumption~2.5, which is assumed in
Proposition~2.3.  Therefore, the import of Proposition~2.4 is:

\medskip
\noindent
If Assumption~2.6 is added to the hypothesis in Proposition~2.3, then the
$\widehat{\mathbf{S}}$ given in \eqref{eq:2.5.1} with
$\mathbf{b} = \widetilde{\bm{\beta}}$ (so $\hat{\varepsilon}_i$ is the OLS residual $e_i$)
can be used in (2.3.5) to calculate the estimated asymptotic variance:
\begin{equation}
  \widehat{\avar}(\mathbf{b}) = \mathbf{S}_{\mathbf{xx}}^{-1}
    \biggl(\frac{1}{n} \sum_{i=1}^{n} e_i^2\, \mathbf{x}_i \mathbf{x}_i'\biggr)
    \mathbf{S}_{\mathbf{xx}}^{-1},
  \label{eq:2.5.4}
\end{equation}
which is consistent for $\avar(\mathbf{b})$.

%% ─── Data Matrix Representation of S ─────────────────────────────────────────
\subsection*{Data Matrix Representation of $\widehat{\mathbf{S}}$}

If $\mathbf{B}$ is the $n \times n$ diagonal matrix whose $i$-th diagonal element is
$\hat{\varepsilon}_i^2$, then the $\widehat{\mathbf{S}}$ in \eqref{eq:2.5.1} can be
represented in terms of data matrices as
\begin{equation}
  \widehat{\mathbf{S}} = \frac{\mathbf{X}' \mathbf{B} \mathbf{X}}{n}
  \tag{2.5.1$'$}
  \label{eq:2.5.1p}
\end{equation}
with
\[
  \mathbf{B} = \begin{bmatrix}
    \hat{\varepsilon}_1^2 & & \\
    & \ddots & \\
    & & \hat{\varepsilon}_n^2
  \end{bmatrix}.
\]
So \eqref{eq:2.5.4} can be rewritten (with $\hat{\varepsilon}_i$ in $\mathbf{B}$ set to
$e_i$) as
\begin{equation}
  \widehat{\avar}(\mathbf{b}) = n \cdot (\mathbf{X}'\mathbf{X})^{-1}
    (\mathbf{X}'\mathbf{B}\mathbf{X})(\mathbf{X}'\mathbf{X})^{-1}.
  \tag{2.5.4$'$}
  \label{eq:2.5.4p}
\end{equation}
These expressions, although useful for some purposes, should not be used for computation
purposes, because the $n \times n$ matrix $\mathbf{B}$ will take up too much of the
computer's working memory, particularly when the sample size is large.  To compute
$\widehat{\mathbf{S}}$ from the sample, the formula \eqref{eq:2.5.1} is more useful than
\eqref{eq:2.5.1p}.

%% ─── Finite-Sample Considerations ────────────────────────────────────────────
\subsection*{Finite-Sample Considerations}

Of course, in finite samples, the power may well be far below one against certain
alternatives.  Also, the probability of rejecting the null when the DGP does satisfy the
null (the Type~I error) may be very different from the assumed significance level.
Davidson and MacKinnon (1993, Section~16.3) report that, at least for the Monte Carlo
simulations they have seen, the robust $t$-ratio based on \eqref{eq:2.5.1} rejects the
null too often and that simply replacing the denominator $n$ in \eqref{eq:2.5.1} by the
degrees of freedom $n - K$ or equivalently multiplying \eqref{eq:2.5.1} by $n/(n - K)$
(this is a \textbf{degrees of freedom correction}) mitigates the problem of overrejecting.
They also report that the robust $t$-ratios based on the following adjustments on
$\widehat{\mathbf{S}}$ perform even better:
\begin{equation}
  \widehat{\mathbf{S}} = \frac{1}{n} \sum_{i=1}^{n}
    \frac{e_i^2}{(1 - p_i)^d}\, \mathbf{x}_i \mathbf{x}_i', \quad d = 1 \text{ or } 2,
  \label{eq:2.5.5}
\end{equation}
where $p_i$ is the $p_i$ defined in the context of the influential analysis in
Chapter~1: it is the $i$-th diagonal element of the projection matrix $\mathbf{P}$,
that is,
\[
  p_i \equiv \mathbf{x}_i'(\mathbf{X}'\mathbf{X})^{-1}\mathbf{x}_i
    = \frac{\mathbf{x}_i' \mathbf{S}_{\mathbf{xx}}^{-1} \mathbf{x}_i}{n}.
\]
